\section{动机：全书回顾与最终目标}
\label{sec:ch10-motivation}
% ============================================================================

走到终章之前，读者其实已经在不知不觉中完成了一次相当长的攀登。
开始时我们面对的只是上半平面上的函数，
以及它们在 $\SL_2(\Z)$ 或 $\Gamma_0(N)$ 作用下的变换性质；
后来这些函数长出 Fourier 展开、Hecke 特征值与 Euler 乘积；
再后来它们又突然变成了模曲线上的微分、Jacobian 上的特征子簇、定义在 $\Q$ 上的代数曲线、
乃至绝对 Galois 群上的二维表示。
如果一路只盯着各章的局部技术，
读者很容易在某一章里感到“这和费马大定理到底还有什么关系”。
终章首先要做的，
就是把这条长链重新从头看到尾。

\textbf{我们遇到了什么问题？}
费马大定理表面上是最古老的整数方程之一：
\[
  x^n+y^n=z^n,
  \qquad n\ge 3.
\]
它的陈述连中学生都能听懂，
可它的证明却躲进了模曲线、Hecke 代数与 Galois 表示之中。
为什么？
因为这个定理的难点不在“某一步代数变形不够巧”，
而在于要识别一种根本不该存在的全局算术对象。
Frey 的洞见正是：
如果费马方程真有解，
它会制造出一条带有极端局部性质的椭圆曲线；
而要判断这条曲线究竟能不能存在，
最有效的办法不是直接算它的点，
而是把它放进模性理论的框架里审判。

\textbf{核心思想是什么？}
前九章已经完成了一次规模很大的“翻译工程”。
第\ref{ch:forms-curves}章到第\ref{ch:hecke}章告诉我们，
模形式空间 $\Sk_k(\Gamma_0(N))$ 里存在由 Hecke 算子同时对角化得到的归一化特征形式 $f$，
它们带着一串极其刚性的系数 $a_p(f)$，也带着 $L$-函数 $L(s,f)$。
第4、5章把这些形式放到模曲线 $X_0(N)$ 上，
从亏格与微分的角度说明尖点形式其实就是几何对象。
第\ref{ch:jacobians}章与第\ref{ch:algebraic-curves}章进一步把 Jacobian $J_0(N)$、模 Abel 簇 $A_f$ 与 $\Q$-模型组织起来，
让“模形式对应几何对象”不再只是类比。
第\ref{ch:eichler-shimura}章证明 Eichler--Shimura 关系，
于是 $L(s,A_f)=L(s,f)^{[K_f:\Q]}$；
第\ref{ch:galois}章再把这件事翻译成
\[
  \rho_{E,\ell}\cong \rho_{f,\ell}
\]
的语言，
说明“椭圆曲线来自模形式”就是 Galois 表示的一致。

\textbf{引入之后能得到什么？}
一旦看清这条主线，
终章的最后一步就显得异常清楚：
只要我们能证明每条半稳定椭圆曲线都模，
那么 Frey 曲线也模；
而一旦它模，Ribet 的水平下降就会把它的 residual 表示一路从级别 $2\,\mathrm{rad}_{\mathrm{odd}}(abc)$ 压到级别 $2$；
但第5章的维数公式已经告诉我们 $\Sk_2(\GammaO(2))=0$。
于是，费马方程若有解，
就会迫使空空间里出现一个新形式。
这不是“算不出来”，而是“概念上不允许存在”。

\begin{figure}[ht]
\centering
\begin{tikzpicture}[>=Stealth, font=\small, node distance=1.4cm]
  \node[draw, rounded corners, fill=blue!8, minimum width=3.3cm, minimum height=0.9cm] (c13) {第1--3章：模形式与 Hecke};
  \node[draw, rounded corners, fill=green!10, right=2.0cm of c13, minimum width=3.3cm, minimum height=0.9cm] (c45) {第4--5章：模曲线与亏格};
  \node[draw, rounded corners, fill=yellow!16, below=1.5cm of c13, minimum width=3.3cm, minimum height=0.9cm] (c67) {第6--7章：Jacobian 与 $\Q$-模型};
  \node[draw, rounded corners, fill=orange!14, right=2.0cm of c67, minimum width=3.3cm, minimum height=0.9cm] (c8) {第8章：Eichler--Shimura 与 $L$-函数};
  \node[draw, rounded corners, fill=purple!12, below=1.5cm of $(c67)!0.5!(c8)$, minimum width=3.6cm, minimum height=0.9cm] (c9) {第9章：Galois 表示与 Ribet};
  \node[draw, rounded corners, fill=red!10, below=1.6cm of c9, minimum width=4.4cm, minimum height=1.0cm] (c10) {第10章：模性定理与费马大定理};

  \draw[->, very thick] (c13) -- node[above] {Fourier 系数} (c45);
  \draw[->, very thick] (c13) -- node[left] {Hecke 特征值} (c67);
  \draw[->, very thick] (c45) -- node[right] {微分与亏格} (c8);
  \draw[->, very thick] (c67) -- node[above] {$A_f$ 与 $J_0(N)$} (c8);
  \draw[->, very thick] (c8) -- node[right] {$L$-函数与 Frobenius} (c9);
  \draw[->, very thick] (c9) -- node[right] {模性与水平下降} (c10);
  \draw[->, thick, dashed] (c67) -- (c9);
  \draw[->, thick, dashed] (c13) to[bend right=18] node[left] {新形式} (c10);
  \draw[->, thick, dashed] (c45) to[bend left=18] node[right] {维数公式} (c10);
\end{tikzpicture}
\caption{全书知识图谱：前九章分别准备了模形式、几何与 Galois 三条主线；终章把它们压缩进同一条证明费马大定理的逻辑链。}
\label{fig:ch10-knowledge-map}
\end{figure}


现在把前九章真正串成一个故事。
第1章的椭圆曲线与模形式给了两个看似平行的世界：
一个世界里有 Weierstrass 方程、挠点与判别式，
另一个世界里有 Fourier 展开、权、级别与尖点。
第3章加入 Hecke 算子之后，
模形式不再只是函数空间里的向量，
而是按特征值系统分解的算术对象。
到此为止，
我们已经拥有“解析数据”。

第4、5章把这个解析世界几何化。
$X_0(N)$ 不只是一个商空间，
而是一条紧黎曼面；
权 $2$ 尖点形式不只是 Fourier 级数，
而是其上的全纯微分。
亏格公式告诉我们，
模形式空间的维数其实是在计算一条曲线的几何复杂度。
这一步非常关键：
它让后来出现的 Jacobian、Abel 映射与模参数化都不再突兀。

第\ref{ch:jacobians}章和第\ref{ch:algebraic-curves}章则把“几何化”推进到算术层面。
Jacobian $J_0(N)$ 把除子类、积分与 Hecke 作用收纳进同一个 Abel 簇中；
与新形式 $f$ 对应的商簇 $A_f$ 则把 Hecke 特征值封装成一个真正的代数几何对象。
与此同时，模曲线与 Hecke 对应在 $\Q$ 上有代数模型，
所以这些对象不仅在复数域上存在，
更能接受 Galois 作用。
这为第\ref{ch:eichler-shimura}章与第\ref{ch:galois}章搭好了桥。

到了第\ref{ch:eichler-shimura}章，
整件事第一次显示出“要走向费马大定理了”的锋芒。
Eichler--Shimura 关系证明，
对权 $2$ 新形式 $f$，模 Abel 簇 $A_f$ 的 Hasse--Weil $L$-函数与 $L(s,f)$ 完全一致。
这意味着模形式的 Fourier 系数不只是解析级数里的数字，
它们还能在有限域点数、Frobenius 特征多项式与椭圆曲线 $L$-函数中重新出现。
而第\ref{ch:galois}章则把这种一致提升到表示论层面：
Deligne 表示与 Tate 模表示满足同样的迹公式，
于是“来自某个模形式”终于能被精确地写成“来自某份二维 Galois 表示”。

这就是终章真正接手的地方。
现在我们不再从零开始，
而是在前九章搭起的舞台上让所有角色同时登场：
费马方程给出 Frey 曲线，
Frey 曲线给出 residual 表示，
Wiles 定理给出模性，
Ribet 定理给出水平下降，
维数公式给出终点空间为空。
整部书的逻辑在这里第一次完全闭环。

\begin{proposition}[终章的证明框架]
\label{prop:ch10-final-roadmap}
设 $p\ge 5$ 为奇素数。
若存在费马方程的非平凡原始解
\[
  a^p+b^p=c^p,
\]
则存在一条半稳定椭圆曲线 $E_{a,b,c}/\Q$ 满足以下性质：
\begin{itemize}
  \item[(i)] $E_{a,b,c}$ 的 odd conductor 为 $\mathrm{rad}_{\mathrm{odd}}(abc)$；
  \item[(ii)] 对每个奇素数 $q\mid abc$，residual 表示 $\bar\rho_{E_{a,b,c},p}$ 在 $q$ 处无分歧；
  \item[(iii)] 若每条半稳定椭圆曲线都模，则 $\bar\rho_{E_{a,b,c},p}$ 来自某个权 $2$ 新形式；
  \item[(iv)] 由水平下降，$\bar\rho_{E_{a,b,c},p}$ 最终来自 $\Sk_2(\GammaO(2))$。
\end{itemize}
由于 $\Sk_2(\GammaO(2))=0$，故不存在这样的费马解。
\end{proposition}

\begin{proof}
(i) 与 (ii) 是下一节 Frey 曲线算术性质的内容；
(iii) 是半稳定模性定理，也就是本章稍后要陈述的 \cref{thm:wiles}；
(iv) 则是第\ref{ch:galois}章证明、并将在本章再次使用的 Ribet 水平下降。
因此，只要这三块拼图同时成立，
费马方程的任何原始解都会制造出一个本不该存在的级别 $2$ 新形式。
空空间不可能产生特征值系统，
所以这样的原始解不存在。
\end{proof}

\textbf{注。}
这个命题本身并不深，
却是整本书最值得反复回望的一页。
它告诉我们：
Wiles 真正需要证明的，
并不是“费马方程没有解”这一句表面陈述，
而是“半稳定椭圆曲线必模”这一条更强、也更结构性的定理。
FLT 由此成为一个更大理论的推论，
而不是那套理论的起点。
